\chapter{用户进程、线程与 libc}
\label{chap:user}

\begin{introduction}[本章内容]
  \item 用户程序的形态与系统调用网关
  \item 从 \code{init} 到 shell
  \item 用户程序一览
  \item 用户线程：\code{clone}、\code{texit}、\code{tjoin}
  \item libc 与 POSIX 源码兼容层
  \item 在 xv6 里编译 C 程序
  \item 一个大型用户程序：\code{sshd}
\end{introduction}

\section{用户程序的形态}

\subsection{为什么需要系统调用}

进程隔离的第一步，是把所有用户程序放到特权最低的 EL0：在这一级，程序不能修改页表基址、不能屏蔽中断、不能访问没有映射给它的内存，也就无法改写内核为隔离所做的任何设置。否则，其他一切隔离手段都没有意义。

可是用户程序总要打印字符、读文件、发网络包，而这些都需要直接操作设备。我们不能让每个程序直接去写 UART 寄存器，而是由内核提供一组受控的入口——系统调用。它不能是普通的函数调用，因为每次调用都要把执行级别从 EL0 提升到 EL1，而提升级别只有一条路：产生异常。所以每个系统调用本质上都是一次\emph{同步异常}：用户程序准备好参数，执行 \code{svc} 指令；硬件把它带到 EL1 的异常向量；内核检查参数、完成请求，然后执行普通的异常返回，从 \code{svc} 的下一条指令继续。进程隔离的第三个方面——每个进程有自己独立的内存视图——由页表完成（第~\ref{chap:mmu} 章）。


\subsection{链接与装载}

xv6 原版把用户程序和用户库\emph{静态}链接成一个 ELF64 文件，链接在虚拟地址 0，入口是 \code{main}。本项目在 \code{make DYNAMIC\_USER=0} 时仍退回这一形态（默认形态改为动态链接，见第~\ref{sec:dynamic-libc} 节）：

\begin{makecode}
_%: %.o $(ULIB)
	$(LD) $(LDFLAGS) -N -e main -Ttext 0 -o $@ $^
\end{makecode}

\code{-N} 让代码段和数据段放在同一个可读写可执行的段里，\code{exec()} 只需处理一个 \code{PT\_LOAD}，也没有 C 运行时启动代码：\code{exec()} 把 \code{argc}、\code{argv} 放进 \code{x0}、\code{x1}，把 PC 设为入口地址，\code{main} 返回后由用户库调用 \code{exit()}——实际上 xv6 的程序习惯在 \code{main} 末尾显式调用 \code{exit(0)}。这一形态简单，但代价是每个程序都各自携带一份 libc 的代码；把用户库改成共享的 \code{libc.so}、让程序默认以位置无关（PIE）\code{ET\_DYN} 形式动态链接它，就是第~\ref{sec:dynamic-libc} 节的内容。

\subsection{系统调用网关}

\file{user/usys.pl} 只生成\emph{一个}私有陷入网关 \code{\_\_syscall}，而不是每个系统调用一条桩。它把系统调用号从第一个参数 \code{x0} 挪进 \code{x7}，其余参数整体下移一格，然后陷入 EL1：

\begin{asmcode}
.global __syscall
__syscall:
 mov x7, x0      // 第 0 个实参是系统调用号
 mov x0, x1      // 真正的参数 x1..x6 下移成 x0..x5
 mov x1, x2
 mov x2, x3
 mov x3, x4
 mov x4, x5
 mov x5, x6
 svc #0          // 陷入 EL1
 ret             // 返回值已在 x0
\end{asmcode}

\file{user/libc\_syscall.c} 里的公开 libc 函数（\code{fork}、\code{read}、\code{write}、\code{exec}……）都是对这个网关的薄包装：一个 \code{sc()} 帮助函数把 \code{SYS\_*} 号作为第一个实参传给 \code{\_\_syscall}，再用 \code{SC0/SC1/\ldots/SC6} 按参数个数展开（\code{I}/\code{P} 只做到 \code{long} 的转换）：

\begin{ccode}
extern long __syscall(long, long, long, long, long, long, long);

static long sc(long number, long a0, long a1, long a2, long a3, long a4, long a5){
  long result = __syscall(number, a0, a1, a2, a3, a4, a5);
  if(result < 0)                 // 内核目前只返回 -1，尚无 Linux 式 -errno
    errno = EIO;
  return result;
}

#define SC3(n,a,b,c) sc(SYS_##n, (a), (b), (c), 0, 0, 0)
int write(int fd, const void *buf, int count)
{ return SC3(write, I(fd), P(buf), I(count)); }
\end{ccode}

这样分层有两个好处。其一，\code{svc} 只出现在这一个网关里，应用和大部分 libc 都不再直接执行陷入指令，系统调用号这套内核 ABI 被收进 \code{libc.so} 内部。其二，在动态链接下（第~\ref{sec:dynamic-libc} 节），应用对 \code{write} 的调用经 PLT 落到共享 libc，再由 libc 内部经 \code{\_\_syscall} 陷入，用户程序镜像里不再各自携带一份桩代码。

内核侧 \code{syscall()} 用 \code{trapframe->x7} 索引函数表，返回值写回 \code{trapframe->x0}。目前共有 70 多个系统调用：除了 xv6 原有的 21 个，还有同步对象、网络（UDP、ICMP、TCP socket、epoll）、挂载、TTY 会话、线程、调度属性、高精度时钟等。

\begin{note}
  系统调用号放在 \code{x7}，这是本项目沿用的约定；Linux AArch64 用的是 \code{x8}。所以 Linux 的二进制或汇编系统调用包装不能直接拿来用，第~\ref{sec:posix} 节的 POSIX 层是\emph{源码}兼容，不是二进制兼容。
\end{note}

\subsection{内核如何分发系统调用}
\label{sec:syscall-dispatch}

\code{svc} 进入 \code{el0sync}，\code{usavereg} 把全部用户寄存器存进 trapframe，\code{usertrap()} 检查 \reg{ESR\_EL1} 的异常类别是否为 \code{0x15}（来自 AArch64 的 SVC），然后打开中断、调用 \code{syscall()}（第~\ref{sec:irq} 节）。分发函数在 \file{kernel/core/syscall.c} 中：

\begin{ccode}
void syscall(void)
{
  int num;
  struct proc *p = myproc();

  num = p->trapframe->x7;
  if(num > 0 && num < NELEM(syscalls) && syscalls[num]) {
    p->trapframe->x0 = syscalls[num]();
  } else {
    printf("%d %s: unknown sys call %d\n", p->pid, p->name, num);
    p->trapframe->x0 = -1;
  }
}
\end{ccode}

\code{syscalls[]} 是一张以系统调用号为下标的函数指针表。处理函数先检查调用号是否越界，越界就返回 $-1$，而不是像教学内核那样直接停机——一个用户程序传错调用号，不应该让整个系统死掉。

参数不需要搬运：用户态执行 \code{svc} 时它们就在 \code{x0}～\code{x5} 中，现在原样躺在 trapframe 里，各个 \code{sys\_*} 函数用 \code{argint()}/\code{argaddr()}/\code{argstr()} 按编号取出：

\begin{ccode}
static uint64 argraw(int n)
{
  struct proc *p = myproc();
  switch (n) {
  case 0: return p->trapframe->x0;
  case 1: return p->trapframe->x1;
  ...
  case 5: return p->trapframe->x5;
  }
  panic("argraw");
  return -1;
}
\end{ccode}

指针参数只取出数值，不在这里检查合法性；真正访问用户内存时由 \code{copyin()}/\code{copyout()} 查用户页表，非法地址会使系统调用返回 $-1$，而不会让内核出错。

返回值写回 \code{p->trapframe->x0}，而不是直接放进 \code{x0} 寄存器。原因是返回用户态时 \code{urestorereg} 会从 trapframe 恢复\emph{所有}寄存器，直接写寄存器的值会被覆盖；只有改写保存的那一份，用户程序在 \code{svc} 之后才能在 \code{x0} 中看到它。

\subsection{两个栈指针与返回级别}

EL0 和 EL1 各有自己的栈指针。从用户态陷入时，硬件立刻把 \code{sp} 换成 \reg{SP\_EL1}（内核栈），用户栈指针原样留在 \reg{SP\_EL0} 中。即使内核从不使用 \reg{SP\_EL0}，也必须把它连同其他寄存器一起保存和恢复：陷入期间可能发生进程切换，下一个返回用户态的进程需要的是\emph{它自己}的用户栈指针。所以 \code{usavereg} 把 \reg{SP\_EL0} 存进 trapframe 的 \code{sp} 字段，\code{urestorereg} 把它写回。

内核态的 \code{savereg} 则不需要恢复 \code{sp}：内核栈是同一个，即使中间发生了进程切换，\code{swtch()} 也已经把 \code{sp} 换成了正确的值。

\code{eret} 回到哪一级，由 \reg{SPSR\_EL1} 的模式字段决定，不需要另外指定，所以总是回到异常发生时所在的那一级。利用这一点，内核可以“伪造”一次从未发生过的异常来进入用户态：\code{userinit()} 在第一个进程的 trapframe 中写入 \code{elr = 0}、\code{spsr = 0}（模式 EL0t）、\code{sp = PGSIZE}，\code{usertrapret()} 一执行 \code{eret}，CPU 就第一次来到 EL0（第~\ref{chap:process} 章）。启动时从 EL2 降到 EL1 用的也是同一个技巧（第~\ref{chap:boot} 章）。

\section{从 \code{init} 到 shell}

第一个用户进程 \code{initcode} 执行 \code{exec("/init")}（第~\ref{chap:process} 章）。\file{user/init.c} 做的事情相当于一个嵌入式 Linux 的启动脚本：

\begin{enumerate}
  \item 创建 \code{/dev}，用 \code{mknod} 建立 \code{console}、\code{tty}、\code{ttyS0}、\code{input/event0}、\code{video0} 设备节点；以 \code{/dev/ttyS0} 打开 fd 0，复制为 fd 1、2；
  \item 创建 \code{/proc}、\code{/boot}、\code{/mnt/ext2}、\code{/etc}、\code{/root}、\code{/bin} 等目录，为根目录下的程序在 \code{/bin} 中建立硬链接；
  \item 首次启动时生成 \code{/etc/hostname}、\code{/etc/fstab}、\code{/etc/passwd}（默认账号 \code{root}，密码 \code{xv6}）、\code{/etc/rc}；
  \item 按 \code{/etc/fstab} 逐行调用 \code{mount()}；
  \item 如果存在 \code{/etc/wifi.conf}，先运行 \code{/bin/wifi} 连接无线网络；
  \item 在一个循环里 fork 并 exec \code{/bin/login}；\code{login} 退出后回收它，再启动一个新的登录提示；同时回收所有过继来的孤儿进程。
\end{enumerate}

\code{login} 校验密码后调用 \code{tty\_attach()} 建立会话，切到 passwd 中指定的家目录，exec 指定的 shell。\code{sh} 在 xv6 原版基础上增加了：读取 \code{/etc/rc}、\code{PATH} 搜索（找不到时再回退到根目录，作为存储异常时的恢复路径）、\code{export}、\code{exit}/\code{logout} 内建命令，以及把外部命令设为前台进程组的作业控制。

\section{用户程序一览}

\begin{table}[htbp]
  \centering
  \caption{用户程序（\code{UPROGS}）}
  \label{tab:uprogs}
  \small
  \begin{tabularx}{\linewidth}{@{}lX@{}}
    \toprule
    类别 & 程序 \\
    \midrule
    基本工具 & \code{sh}、\code{init}、\code{login}、\code{ls}、\code{cat}、\code{echo}、\code{grep}、\code{wc}、\code{mkdir}、\code{rm}、\code{ln}、\code{mv}、\code{touch}、\code{kill}、\code{ps}、\code{edit}、\code{file} \\
    文件系统 & \code{ext2ls}、\code{ext2cat}、\code{stressfs} \\
    网络 & \code{wifi}、\code{dhcp}、\code{ping}、\code{netdns}、\code{nettest}、\code{tftp}、\code{tcpd}、\code{epollserver}、\code{sshd} \\
    调度与同步 & \code{chrt}、\code{taskset}、\code{nice}、\code{renice}、\code{rtlat}、\code{spin}、\code{syncdemo}、\code{prodcons} \\
    线程与 POSIX & \code{threadtest}、\code{posixtest} \\
    设备 & \code{camshot}、\code{ptytest}、\code{vmmap} \\
    开发 & \code{tcc} \\
    测试 & \code{usertests}、\code{forktest}、\code{grind}、\code{zombie} \\
    \bottomrule
  \end{tabularx}
\end{table}

几个调度工具的用法与 Linux 相同，它们都是“先修改自己的调度属性，再 exec 目标程序”，利用了 \code{fork}/\code{exec} 对调度属性的继承：

\begin{shcode}
chrt -f 80 rtlat          # 以 SCHED_FIFO 优先级 80 运行 rtlat
taskset 0x8 spin &        # 只允许在 CPU3 上运行
nice -n 10 grind          # 降低普通进程的时间片
chrt -f -p 60 6           # 修改 PID 6（某个 kworker）的策略
\end{shcode}

\section{用户线程}

\subsection{设计}

原来的模型是“一个 \code{struct proc} 对应一个独立的地址空间”。加入线程时，项目没有引入新的“线程控制块”，而是让每个线程仍然是一个完整的 \code{struct proc}——有自己的内核栈、trapframe、调度属性、TID 和 CPU 亲和性——只是同一线程组的成员共享同一个引用计数的 \code{struct vmspace}（图~\ref{fig:threads}）。

\begin{figure}[htbp]
  \centering
  \begin{tikzpicture}[node distance=4mm and 14mm]
    \node[box, text width=30mm] (p) {leader\\TID $=$ TGID};
    \node[box, below=of p, text width=30mm] (t1) {线程 1\\独立内核栈/trapframe};
    \node[box, below=of t1, text width=30mm] (t2) {线程 2\\独立内核栈/trapframe};
    \node[hbox, right=of t1, text width=36mm] (vm) {\code{struct vmspace}\\TTBR0 页表、sz\\refcount $=3$};
    \draw[arr] (p.east) -- (vm); \draw[arr] (t1) -- (vm); \draw[arr] (t2.east) -- (vm);
  \end{tikzpicture}
  \caption{线程组共享 \code{vmspace}}
  \label{fig:threads}
\end{figure}

调度器切换到任何一个线程时，\code{switchuvm()} 都写入 \code{p->vm->pagetable}，所以同组线程看到相同的代码、全局变量、堆和彼此的栈；而它们的内核栈和 trapframe 各自独立，可以同时在多个 Cortex-A53 核上运行。

\subsection{系统调用}

\begin{table}[htbp]
  \centering
  \caption{线程相关系统调用}
  \small
  \begin{tabularx}{\linewidth}{@{}lX@{}}
    \toprule
    调用 & 作用 \\
    \midrule
    \code{clone(entry, arg, stack\_top)} & 创建共享地址空间的新调度实体：\code{x0}=arg，\code{SP}=stack\_top，\code{ELR}=entry \\
    \code{texit(status)} & 只退出调用线程；leader 调用时等同于整个进程退出 \\
    \code{tjoin(tid, \&status)} & 等待并回收同组中的指定线程（0 表示任意） \\
    \code{gettid()} / \code{getpid()} & 返回线程 ID / 线程组 ID \\
    \bottomrule
  \end{tabularx}
\end{table}

用户库在其上提供更安全的包装，它负责分配 16 字节对齐的 16\,KiB 用户栈，并保证线程函数返回后自动调用 \code{texit()}：

\begin{ccode}
static void thread_boot(void *opaque) {
  struct thread_start start = *(struct thread_start *)opaque;
  free(opaque);                       // 启动记录用完即释放
  start.fn(start.arg);
  texit(0);                           // 线程函数返回后自动退出
}

int thread_create(thread_t *thread, void (*fn)(void *), void *arg) {
  thread->stack = malloc(THREAD_STACK_SIZE + 16);
  struct thread_start *start = malloc(sizeof(*start));
  start->fn = fn;  start->arg = arg;
  uint64 top = ((uint64)thread->stack + THREAD_STACK_SIZE + 16) & ~15ULL;
  thread->tid = clone(thread_boot, start, (void *)top);
  ...
}
\end{ccode}

用户栈要等 \code{thread\_join()} 中内核回收了线程之后才释放，否则线程可能还在用这块栈。

\subsection{\code{clone} 为什么“返回两次”}

\code{fork()} 会返回两次：父进程中返回子进程的 pid，子进程中返回 0，因为子进程的 trapframe 是父进程的完整拷贝，只把 \code{x0} 改成了 0。Linux 的 \code{clone} 系统调用语义相同，所以 glibc 的 \code{clone()} 包装函数要在汇编里处理两次返回：调用前把函数指针和参数存进不会被系统调用破坏的寄存器；系统调用返回非 0 说明还在原线程，直接返回调用者；返回 0 说明已经身处新线程，就取出函数指针调用它，函数结束后再调用 \code{exit}——新线程永远不会回到包装函数的调用者那里。

xv6-aarch64 的 \code{clone(entry, arg, stack\_top)} 换了一种更简单的语义：内核直接把新线程的入口、参数和栈写进它的 trapframe，新线程第一次回到用户态时就从 \code{entry} 开始执行，用户态不需要任何汇编包装：

\begin{ccode}
int threadclone(uint64 entry, uint64 arg, uint64 stack_top)
{
  ...
  if(entry == 0 || stack_top < 16 || (stack_top & 15) != 0)
    return -1;                          // 栈必须 16 字节对齐
  ...
  vm->refcount++;                       // 共享同一个地址空间
  vmspace_put(np->vm);
  np->vm = vm;
  np->is_thread = 1;
  np->tgid = p->tgid;

  *(np->trapframe) = *(p->trapframe);
  np->trapframe->elr = entry;           // eret 之后从这里开始
  np->trapframe->sp  = stack_top;       // 新的用户栈
  np->trapframe->x0  = arg;             // 第一个参数
  np->trapframe->x30 = 0;               // 不允许“返回”
  ...
  make_runnable(np, 0);
  return tid;
}
\end{ccode}

\code{x30} 被清零，线程函数如果直接 \code{ret} 就会跳到地址 0 出错，所以用户库的 \code{thread\_boot()} 包了一层，保证函数返回后调用 \code{texit()}。线程退出后变成 ZOMBIE，等 \code{tjoin()} 回收——与进程的 \code{exit}/\code{wait} 是同一种约定：先保留“尸体”，让等待者能取到退出状态，再由等待者释放资源。

\subsection{多核下的细节}

\begin{itemize}
  \item \code{sbrk()} 修改共享的页表时持有 \code{vmspace.lock}，改完执行 \code{tlbi vmalle1is} 广播失效，覆盖其他核上正在运行的同组线程；
  \item 用户态的 K\&R \code{malloc} 加了一把基于 \code{ldaxr}/\code{stxr} 的锁；
  \item \code{exec()} 只在线程组只剩调用者时才允许执行；\code{exit()} 先终止并回收同组线程；\code{kill(pid)} 以 TGID 为目标，标记并唤醒整组；
  \item \code{clone()} 像 \code{fork()} 一样复制文件描述符\emph{引用}——底层 \code{struct file} 和偏移共享，但描述符表独立。这与 Linux 的 \code{CLONE\_FILES} 不同，是有意保留的简化。
\end{itemize}

\code{threadtest} 创建两个线程，用内核同步对象保护共享计数器，允许它们被调度到不同 CPU，预期输出 \code{threadtest: PASS counter=2000}。

\subsection{实例：threadtest 的运行过程}
\label{sec:threadtest-fig}

\begin{figure}[p]
  \centering
  \includegraphics[height=0.93\textheight,keepaspectratio]{threadtest.pdf}
  \caption{threadtest：exec 之后的两份快照（pid 20），以及 clone 出两个线程之后三个 proc 共享一个 vmspace 的结构}
  \label{fig:threadtest-run}
\end{figure}

图~\ref{fig:threadtest-run} 的上半部分是 \code{exec("/bin/threadtest")} 时打印的两份快照，与第~\ref{sec:fork-snapshot} 节 ls 的情况几乎相同。struct proc、vmspace、内核栈用的正是 ls 退出后释放的那几页。有一点不同：新页表的物理页这次是递增分配的（\code{0x6e99} $\to$ \code{0x6ea4}），因为 ls 退出时按顺序还回了页，kalloc 后进先出，取出来的顺序就反过来了。

下半部分是按 \code{thread.c} 与 \code{threadclone()} 推出的线程结构，线程 tid 和线程栈地址为示意。\code{thread\_create()} 先用 malloc 在堆里分配 16\,KiB 的线程栈和一个 start 记录，再调用 \code{clone(thread\_boot, start, 栈顶)}。内核为新线程分配独立的 struct proc、内核栈、trapframe 和 context，但让 \code{np->vm} 指向同一个 vmspace（\code{refcount} 从 1 增到 3），因此三个线程共用页表根 \code{0x6e99000}。新线程的 trapframe 从主线程复制后被改写：\code{elr = thread\_boot}（libc 中 \code{0xe35c}），\code{sp} 设为线程栈顶，\code{x0} 设为 start 记录。线程第一次被调度时经 \code{forkret} $\to$ \code{usertrapret} $\to$ \code{eret} 进入 \code{thread\_boot}，再进入 \code{worker}。全局变量 \code{counter} 和 \code{counter\_lock} 位于程序的 .bss（\code{0x1a04}、\code{0x1a00}），所有线程共享，所以要用 mutex 保护。线程栈只是堆中的普通内存，下方没有保护页，溢出时会悄悄覆盖相邻的堆数据。

\section{libc 与 POSIX 源码兼容层}
\label{sec:posix}

\subsection{\code{libc.a}}

在最简形态里，用户库被打包成一个静态库 \code{libc.a}（也就是 \code{make DYNAMIC\_USER=0} 时保留的回退路径）：

\begin{itemize}
  \item \file{ulib.c}：字符串、内存函数、\code{gets}、\code{stat}、\code{atoi}；
  \item \file{usys.S}：唯一的 \code{\_\_syscall} 陷入网关；
  \item \file{printf.c}：\code{printf}/\code{fprintf}，支持 \code{\%d \%x \%p \%s \%c}；
  \item \file{umalloc.c}：K\&R 风格的 \code{malloc}/\code{free}，用 \code{sbrk} 向内核要内存；
  \item \file{thread.c}：上一节的线程包装；
  \item \file{posix.c}：POSIX 包装层；
  \item \file{tweetnacl.c}、\file{ssh\_crypto.c}：密码学原语（\code{sshd} 使用）。
\end{itemize}

\subsection{动态链接 \code{libc.so}}
\label{sec:dynamic-libc}

上面的静态 \code{libc.a} 让每个程序各自携带一份 libc 代码。本项目把默认形态改成共享库：\file{Makefile} 先把每个 libc 源文件编成位置无关对象（\code{-fPIC}），再用 \code{ld -shared} 链成 \code{ET\_DYN} 的 \code{libc.so}，SONAME 取名 \code{libc.so}；用户程序则以 \code{-pie} 链接到它：

\begin{makecode}
$U/%.pic.o: $U/%.c
	$(CC) $(CFLAGS) -fPIC -c -o $@ $<
$U/libc.so: $(LIBC_PIC_OBJS)
	$(LD) $(LDFLAGS) -shared --no-undefined -soname libc.so -o $@ $^
_%: %.o $(ULIB) $(USER_LINK_STAMP)
	$(LD) $(LDFLAGS) -pie --no-dynamic-linker -e main -o $@ ...
\end{makecode}

关键在最后一条：\code{-pie} 产出 \code{ET\_DYN} 可执行文件，\code{--no-dynamic-linker} 则\emph{不}写 \code{PT\_INTERP}——因为本项目没有用户态 \code{ld.so}，动态装载的活由内核 \code{exec()} 自己完成。要退回上一节的静态形态，只需 \code{make DYNAMIC\_USER=0}；构建用 \file{.dynamic-user}/\file{.static-user} 两个模式戳保证切换时重链所有用户程序，根文件系统构建再把共享库安装到 \file{/lib/libc.so}。

\paragraph{\code{exec()} 里的装载流程.} 对每个 \code{ET\_DYN} 程序，\file{kernel/fs/exec.c}：

\begin{enumerate}
  \item 校验 ELF 头后按程序头把主镜像映射到虚拟地址 0，并记下 \code{PT\_DYNAMIC} 位置。为此 \code{loadseg()} 放宽了老 xv6 的页对齐要求，只校验 GNU \code{ld} 保证的 \code{p\_vaddr $\equiv$ p\_offset (mod p\_align)}，以容纳非页对齐的段边界。
  \item \code{parsedynamic()} 逐条读 \code{DT\_*} 项，收集 \code{DT\_NEEDED}、\code{DT\_SYMTAB}、\code{DT\_STRTAB}、\code{DT\_HASH}、\code{DT\_RELA/RELASZ/RELAENT} 与 \code{DT\_JMPREL/PLTRELSZ/PLTREL}；符号数取自 ELF hash 表头，动态表必须以 \code{DT\_NULL} 结束。
  \item 策略校验：\code{DT\_NEEDED} 必须恰好一个、且字符串正是 \code{"libc.so"}，随后打开硬编码路径 \file{/lib/libc.so}。任意依赖搜索留给未来的用户态 \code{ld.so}，不藏进内核。
  \item libc 的载入基址取 \code{libbase = PGROUNDUP(sz)}，即紧跟主镜像之后的第一个空闲页，\code{loadimage()} 据此把 \code{libc.so} 也映射进本页表。
  \item \code{relocateimage()} \emph{先}重定位 libc、\emph{再}重定位主镜像：处理 \code{.rela.dyn} 与 \code{.rela.plt}，符号解析顺序为“主镜像优先、libc 兜底”，未定义的强符号会让 \code{exec()} 直接失败。
  \item 校验入口合法、布置 argv 栈、把 PC 设为 \code{mainimg.base + e\_entry}；由于重定位用 \code{copyout} 改写了用户代码与 GOT 页，返回用户态前必须 \code{uvmsync\_icache()} 做 I-cache 维护。
\end{enumerate}

\paragraph{支持的重定位类型.} 装载器只认这份 AArch64 libc 实际会生成的四种 \code{R\_AARCH64\_*}，且全部在装载期一次性（eager）绑定，没有延迟 PLT：

\begin{ccode}
switch(type){
case R_AARCH64_NONE:                       // 跳过
  continue;
case R_AARCH64_RELATIVE:                   // 基址 + addend
  value = img->base + rela.addend; break;
case R_AARCH64_ABS64:                      // 解析符号后 + addend
case R_AARCH64_GLOB_DAT:
case R_AARCH64_JUMP_SLOT:
  value = resolvesym(...) + rela.addend; break;
default:                                   // 其它类型一律失败
  printf("exec: unsupported AArch64 dynamic relocation %d\n", type);
  return -1;
}
\end{ccode}

真机上以默认配置启动、登录后，串口打印每个动态程序里 \code{libc.so} 的载入基址，这两行即证明主 ELF 与 \code{libc.so} 都被成功装载、符号被解析、重定位已完成：

\begin{txtcode}
exec: dynamic /bin/login, libc.so base=0x0000000000002000
xv6-rpi3 login: root
Password: xv6
Welcome root
exec: dynamic /bin/sh, libc.so base=0x0000000000003000
\end{txtcode}

这里的 \code{base} 就是上文的 \code{libbase}：\code{login} 主镜像最大到 \code{0x1a20}，页对齐后落在第 2 页 \code{0x2000}；\code{sh} 更大，跨到第 3 页 \code{0x3000}。若某个程序没有打印 \code{exec: dynamic} 却能运行，说明它是静态 \code{ET\_EXEC} 回退路径。

\paragraph{收敛与限制.} 这是一个刻意为之的最小 \textbf{xv6 原生}装载 ABI，而不是完整的 \code{ld.so}：只支持 RELA（无 \code{DT\_REL}）；只认上面四种动态重定位，带 GOT 数据重定位（\code{ADR\_GOT\_PAGE}/\code{LD64\_GOT\_LO12\_NC}）的二进制会失败；只允许受信任的单个 \file{/lib/libc.so}，libc 自身不得再有依赖；没有 \code{PT\_INTERP}、\code{dlopen/dlsym}、延迟 PLT 绑定、\code{IFUNC}、TLS、符号 interposition 与 ASLR；\code{libc.so} 每次 \code{exec} 都重新读盘并复制进本进程地址空间，进程间不共享只读物理页。于是它消除了每个程序静态携带 libc 的磁盘冗余，却还没省下运行时的物理内存。可插拔 \code{.ko} 模块的装载器与它共用 \file{kernel/elf.h} 里的数据结构，却是两套独立逻辑：模块装载器处理 \code{ET\_REL}、做的是指令位域修补，而这里处理 \code{ET\_DYN}、做的是把基址与符号值写回 64 位槽。

\subsection{一次库函数调用：PLT 与 GOT}
\label{sec:plt-got}

动态链接最核心的问题是：程序编译时不知道 \code{strlen} 会被装到哪里，代码里又不能留绝对地址（否则代码页就要按进程改写，无法共享）。AArch64 的答案是“经过一张表间接跳转”。以 \file{/bin/echo} 为例，先看 ELF 里记了什么：

\begin{txtcode}
$ readelf -lW user/_echo
  PHDR     vaddr 0x40     0x188  R      程序头表自身
  INTERP   vaddr 0x1c8    0x00b  R      "/lib/ld.so"
  LOAD     vaddr 0x0      0x3ec  R E    ELF 头、.plt、.text、只读数据
  LOAD     vaddr 0x1ec0   0x158  RW     .dynamic、.got、.got.plt
$ readelf -dW user/_echo
  NEEDED  libc.so    HASH 0x1d8   SYMTAB 0x200   STRTAB 0x260
  PLTGOT  0x1fe8     JMPREL 0x280 PLTRELSZ 72    PLTREL RELA
$ readelf -rW user/_echo
  0x2000  R_AARCH64_JUMP_SLOT  write
  0x2008  R_AARCH64_JUMP_SLOT  exit
  0x2010  R_AARCH64_JUMP_SLOT  strlen
\end{txtcode}

\code{main} 调用 \code{strlen} 时，\code{bl} 的目标是本程序里一段 16 字节的桩代码——PLT 项；桩代码从 GOT 读出真实地址再跳过去：

\begin{asmcode}
 368:  bl   310 <strlen@plt>       // main 里：调用 PLT 项
0000000000000310 <strlen@plt>:
 310:  adrp x16, 2000              // GOT 所在页，PC 相对
 314:  ldr  x17, [x16, #16]        // x17 = GOT[0x2010]
 318:  add  x16, x16, #0x10
 31c:  br   x17                    // 跳进 libc
\end{asmcode}

GOT 在可写段里，构建出来时三个槽都指向 \code{0x2d0}，也就是 PLT0——glibc 做\emph{延迟绑定}时的解析入口。本项目不做延迟绑定，PLT0 依赖的保留槽是 0，真跳过去会崩溃，所以必须在 \code{main} 之前把槽写好（急切绑定）。处理 \code{R\_AARCH64\_JUMP\_SLOT strlen @ 0x2010} 时，在全局作用域里查到 libc 的 \code{strlen}（\code{st\_value} $=$ \code{0x4260}），libc 若被映射在 \code{0x10\_0000\_1000}，就写入 \code{0x10\_0000\_5260}。此后每次库调用多一次 \code{ldr} 和一次 \code{br}，换来的是代码里没有任何绝对地址。

libc 自己也一样。它以 \code{-shared}、不带 \code{-Bsymbolic} 链接，内部对导出函数的调用同样走自己的 PLT（例如 \code{putc} 里的 \code{bl write@plt}），因此有 71 条 \code{JUMP\_SLOT}，目标全在 RW 段的 \code{.got.plt}。这样程序可以覆盖 libc 的函数：符号查找先找程序，程序若自己定义了 \code{malloc}，libc 内部的 \code{malloc} 调用也会解析到它。表~\ref{tab:dynrel} 是装载器支持的全部重定位类型。

\begin{table}[htbp]
  \centering
  \caption{支持的动态重定位（S：符号地址，A：addend，B：模块基址）}
  \label{tab:dynrel}
  \small
  \begin{tabularx}{\linewidth}{@{}lllX@{}}
    \toprule
    类型 & 编号 & 计算 & 典型位置 \\
    \midrule
    \code{R\_AARCH64\_NONE} & 0 & 跳过 & \\
    \code{R\_AARCH64\_ABS64} & 257 & S $+$ A & 数据里的绝对指针 \\
    \code{R\_AARCH64\_GLOB\_DAT} & 1025 & S $+$ A & \code{.got}：全局变量地址 \\
    \code{R\_AARCH64\_JUMP\_SLOT} & 1026 & S $+$ A & \code{.got.plt}：函数地址 \\
    \code{R\_AARCH64\_RELATIVE} & 1027 & B $+$ A & 模块内部的指针 \\
    \bottomrule
  \end{tabularx}
\end{table}

\subsection{共享 libc 代码页}

上面的装载器让每个进程都有一份私有的 \code{libc.so}。可是 libc 是位置无关的、没有 \code{DT\_TEXTREL}：所有重定位都落在它的 RW 段（GOT 与 \code{.data}），只读页在任何基址下都逐字节相同。用 \code{readelf} 看这份 libc：

\begin{txtcode}
LOAD  vaddr 0x0     memsz 0x9272   R-X   代码与只读数据：第 0..9 页
LOAD  vaddr 0xa278  memsz 0x390    RW-   GOT、.data、.bss：落在 0xa000 这一页
重定位：71 条 R_AARCH64_JUMP_SLOT，目标全部在 RW 段
\end{txtcode}

所以只读页可以让所有进程共用一份物理内存，RW 页仍每个进程一份。第一版实现在 \code{exec()} 里专门缓存 \file{/lib/libc.so} 的代码页，第二次以后直接映射；现在这件事交给通用的页缓存（第~\ref{sec:sharedpages} 节）：\code{shlib\_map()} 对每个代码页调用 \code{pcache\_getpage()}，以用户只读、可执行、PXN 映射，并用 \code{shlib\_textsize()} 保证只读段满足 \code{p\_vaddr == p\_offset}、且不与可写页共用一页。重定位若要写进共享的代码区，\code{exec()} 直接拒绝——那会让“共享”的页带上某个进程的基址。

（当前源码树改用用户态解释器 \file{ld-xv6.so} 之后，这段内核代码已经删除，libc 改由解释器加载，见第~\ref{sec:ldxv6} 节。）

真机上 \code{login} 和 \code{sh} 的 libc 基址不同（\code{0x2000} 与 \code{0x3000}，由各自主镜像的大小决定），这并不妨碍共享：代码页内容与装载地址无关，只是各进程把同一组物理页映射到不同的虚拟地址。

\subsection{用户态解释器 \file{/lib/ld-xv6.so}}
\label{sec:ldxv6}

内核里的装载器只认一个 \file{/lib/libc.so}，还把 ELF 解析、符号查找、重定位都放在内核里。Linux 的做法是交给用户态：程序头里的 \code{PT\_INTERP} 指名一个解释器，内核只映射主程序和解释器，其余由解释器在 EL0 完成。当前源码树就是这样做的，解释器叫 \file{/lib/ld-xv6.so}（\file{user/ldso.c}，约 350 行）。

\paragraph{构建.} 程序以 \code{-pie -dynamic-linker /lib/ld-xv6.so} 链接，于是带上 \code{PT\_INTERP}。解释器由 \file{ldso.pic.o} 和 \file{usys.pic.o} 链成，只借用 libc 的 \code{\_\_syscall} 网关：

\begin{makecode}
$U/ld-xv6.so: $U/ldso.pic.o $U/usys.pic.o
	$(LD) $(LDFLAGS) -shared -Bsymbolic -e ldso_start -o $@ $^
\end{makecode}

\code{mkfs} 新增 \code{-I} 选项，把它放进 \file{/lib}；ext2 根镜像也一并写入。\code{readelf} 显示它只有两个 PT\_LOAD，入口 \code{0x750} 就是 C 函数 \code{ldso\_start}，而且\emph{没有任何重定位}：\code{-fPIC} 加 \code{-Bsymbolic} 让所有自身引用都在链接时绑定为 PC 相对访问。因此内核把它装到任何基址都能直接运行，它也不必像 glibc 的 \code{ld.so} 那样先重定位自己。

\paragraph{内核这一半.} \code{exec()} 删掉了约 460 行内核装载代码，只剩三件事：

\begin{enumerate}
  \item \code{readinterp()} 找到唯一的 \code{PT\_INTERP}，路径以 NUL 结尾、不超过 128 字节；\code{ET\_DYN} 程序没有它就拒绝（\code{exec: dynamic \ldots has no valid PT\_INTERP}）；
  \item 主程序装到 VA 0，解释器装在 \code{interpbase = PGROUNDUP(sz)}，两者都是 \code{loadimage()} 分配的私有页；入口改为 \code{interpbase + e\_entry}；
  \item 初始栈在 argv 之后放空的 envp 和 Linux 形式的辅助向量，argc 仍由 \code{x0} 传递：
\end{enumerate}

\begin{txtcode}
sp -> argv[0..argc-1], NULL        <- x1；x0 = argc
      envp[0] = NULL
      AT_PHDR = e_phoff  AT_PHENT = 56  AT_PHNUM  AT_PAGESZ = 4096
      AT_BASE = interpbase（未使用）     AT_ENTRY = main   AT_NULL
\end{txtcode}

\paragraph{解释器这一半.} \code{ldso\_start(argc, argv)} 用 \code{argv + argc + 2} 找到 auxv，然后：

\begin{enumerate}
  \item 由 \code{AT\_PHDR} 找到主程序的 \code{PT\_DYNAMIC}（主程序基址视为 0），解析 \code{DT\_*}，符号个数取自 \code{DT\_HASH} 的 \code{nchain}；
  \item 要求 \code{DT\_NEEDED} 恰好一个且就是 \code{"libc.so"}；
  \item \code{load\_libc()}：\code{open("/lib/libc.so")}，用 \code{pread} 读 ELF 头和程序头，按最大段尾 \code{0xa650} 调用 \code{sbrk(0xb000)} 向堆要 11 页，再把每个 PT\_LOAD 用 \code{pread} 读到 \code{base + p\_vaddr}；\code{sbrk} 的新页已经清零，\code{.bss} 不用另外处理；
  \item 先重定位 libc（71 条 \code{JUMP\_SLOT}），再重定位主程序；符号在“主程序 $\to$ libc”中\emph{线性}查找，因此程序可以覆盖 libc 的函数；
  \item 向 stderr 打印 \code{ld.so: relocated libc.so in user space}，以 \code{(argc, argv)} 调用 \code{AT\_ENTRY}，再用返回值 \code{exit}。
\end{enumerate}

按这次构建推算，\file{/bin/sh} 的地址空间是：程序 \code{0x0}--\code{0x2e10}，\file{ld-xv6.so} \code{0x3000}--\code{0x4d20}，保护页和 4 页栈 \code{0x5000}--\code{0x9fff}，libc \code{0xa000}--\code{0x14fff}，\code{malloc} 的堆从 \code{0x15000} 开始。所有页都是这个进程私有的，可读、可写、可执行。

\begin{table}[htbp]
  \centering
  \caption{三种 libc 装载方式}
  \label{tab:ldso-compare}
  \small
  \begin{tabularx}{\linewidth}{@{}lXXX@{}}
    \toprule
    & 内核加载（第一步） & \file{ld-xv6.so}（当前） & \file{ld.so}（实验分支） \\
    \midrule
    谁加载 & 内核 \code{exec} & 用户态，\code{pread} & 用户态，\code{mmap} \\
    libc 位置 & 主程序之后 & \code{sbrk} 堆起点 & mmap 区 \\
    代码页 & 跨进程共享 & 每进程一份 & 页缓存共享 \\
    fork & 整页复制 & 整页复制 & 写时复制 \\
    依赖 & 仅 libc & 仅 libc & 多个，\file{/lib} 搜索 \\
    内核需要 & 装载器 & \code{pread}、\code{sbrk} & mmap、页缓存、COW \\
    \bottomrule
  \end{tabularx}
\end{table}

\paragraph{问题与建议.} 对内核要求最少是 \file{ld-xv6.so} 最大的优点，表~\ref{tab:ldso-compare} 也列出了它的代价。分析代码时还发现以下几点：

\begin{itemize}
  \item \textbf{I-cache 一致性（真机隐患）}：libc 的代码是由内核 \code{copyout} 经数据缓存写进堆页的，之后没有做 I-cache 维护——\code{exec} 末尾的 \code{uvmsync\_icache()} 在解释器运行之前就已结束，而 \reg{SCTLR\_EL1} 没有设置 UCI（bit 26），EL0 不能自己执行 \code{dc cvau}/\code{ic ivau}。Cortex-A53 的 I-cache 不侦听 D-cache，取指可能拿到旧内容，正是第~\ref{sec:caches} 节故障二那一类问题；QEMU 不模拟缓存，所以看不出来。修法：像 Linux 一样打开 UCI，让解释器在读完代码段后逐行 \code{dc cvau}、\code{dsb ish}、\code{ic ivau}、\code{dsb ish; isb}；或者由内核提供一个 \code{cacheflush(addr, len)} 系统调用。
  \item \textbf{libc 与堆相邻}：用户 \code{sbrk} 负数缩到 libc 以下会把 libc 一起释放；每个进程多占 11 页，fork 时还要整页复制。
  \item \textbf{依赖“无重定位”}：以后若在 \file{ldso.c} 里写出带指针初值的全局变量，链接器会生成 \code{R\_AARCH64\_RELATIVE}，而解释器不处理它，结果是\emph{静默}出错。建议在 Makefile 里检查 \code{readelf -r} 为空，或者用 \code{AT\_BASE} 做自重定位。
  \item \textbf{模式戳未改名}：\code{exec} 现在拒绝没有 \code{PT\_INTERP} 的 \code{ET\_DYN}，但 \file{.dynamic-user} 戳没变；从旧链接方式切换过来时如果没有 \code{make clean}，开机会报 \code{exec: dynamic /init has no valid PT\_INTERP}，随后 \code{panic: init exiting}。
  \item 每次 \code{exec} 都会打印两行（内核的 \code{-> interpreter} 和解释器的 \code{relocated libc.so}）；\code{AT\_PHDR} 用 \code{e\_phoff} 而不是 \code{PT\_PHDR}；报错前缀仍是 \code{ld.so:}。
\end{itemize}

\subsection{图解：\file{/bin/echo} 的文件布局与运行时绑定}
\label{sec:echo-figs}

本小节用两张图把前面的内容串起来。图中的偏移、地址和符号值都取自这次构建的 \code{readelf -h -l -S -d -r --dyn-syms} 和 \code{objdump -DD user/\_echo}；libc 基址 \code{0x9000} 是按下文推算的值。

\begin{figure}[p]
  \centering
  \includegraphics[height=0.93\textheight,keepaspectratio]{echo-elf-detail.pdf}
  \caption{\file{\_echo} 的文件布局：ELF 头、程序头（段）、节、节头，以及它们之间的指针}
  \label{fig:echo-layout}
\end{figure}

\paragraph{文件布局.} 图~\ref{fig:echo-layout} 的中间一列是文件内容，按偏移从上到下排列（未按比例），每块都写出了解码后的内容。\code{.comment} 和 9 个 \code{.debug\_*} 节（\code{-g} 生成，偏移 \code{0x4b0}--\code{0xa3a}）与加载无关，图中省略。

左侧箭头表示三张头部表指向文件的哪个位置：
\begin{itemize}
  \item ELF 头的 \code{e\_entry = 0x2c8} 指向 \code{main}，\code{e\_phoff = 0x40} 指向程序头表，\code{e\_shoff = 0xf98} 指向节头表；
  \item 程序头（橙色）给 \code{exec} 和 ld-xv6 看，按\emph{段}描述；
  \item 节头（蓝色）给 \code{ld}、\code{readelf}、\code{gdb} 看，按\emph{节}描述。
\end{itemize}

右侧箭头表示节与节之间的引用：
\begin{itemize}
  \item 紫色：\code{bl} 跳到 PLT，\code{adrp/add} 取 \code{.rodata} 中的字符串，这些在链接时已经算好；
  \item 红色：PLT 存根读取 GOT 槽，GOT 槽的初值 \code{0x278} 又指回 PLT0；
  \item 青色：\code{.rela.plt} 的 \code{r\_offset} 指向 GOT 槽，\code{sym} 指向 \code{.dynsym}，\code{.dynsym} 的 \code{st\_name} 指向 \code{.dynstr}；
  \item 绿色虚线：\code{.dynamic} 的 \code{\{tag, val\}} 指向各张表。
\end{itemize}

最右两条色带分别标出两个 LOAD 段的覆盖范围，以及 PHDR、INTERP、DYNAMIC 三个“指针段”指向的内容。从图中可以读出以下几点：

\begin{enumerate}
  \item \textbf{加载只看程序头。} \code{exec} 遍历程序头表，只处理 \code{PT\_LOAD}：用 \code{vaddr + memsz} 调用 \code{uvmalloc}，用 \code{(vaddr, off, filesz)} 调用 \code{loadseg}。节头表、\code{.symtab}、\code{.strtab} 不在任何 LOAD 段里，\code{strip} 掉以后程序照样运行。
  \item \textbf{节被打包进段。} \code{.interp} 到 \code{.rodata} 共 8 个节合成 LOAD~\#1（R-X，文件 \code{0x000}--\code{0x358} 映射到 VA 0）。\code{.dynamic}、\code{.got}、\code{.got.plt} 合成 LOAD~\#2（RW-，映射到 VA \code{0x1358}）。链接时没有用 \code{-z separate-code}，所以只读元数据、代码和常量同在一个 R-X 段。
  \item \textbf{第二个 LOAD 段的 VA 不按页对齐。} ELF 只要求 \code{(vaddr - off) \% align == 0}（\code{0x1358 - 0x358 = 0x1000}），所以 \code{exec} 去掉了旧版的 \code{vaddr \% PGSIZE} 检查，并让 \code{loadseg} 能从页中间开始写入。\code{\_echo} 的 \code{FileSiz == MemSiz}，没有 \code{.bss}。
  \item \textbf{INTERP、DYNAMIC、PHDR 不单独加载。} 它们都落在某个 LOAD 段内部，只起指路作用：INTERP 告诉内核先运行 \file{/lib/ld-xv6.so}；PHDR 和 \code{AT\_PHDR} 相减得到主程序基址；DYNAMIC 是 ld-xv6 的目录。运行时没有节头可用，ld-xv6 只能通过 \code{DT\_SYMTAB}、\code{DT\_HASH}、\code{DT\_JMPREL} 等找到各张表。
  \item \textbf{PLT0 和 \code{GOT[0..2]} 用不上。} 它们是为惰性绑定保留的。ld-xv6 在程序开始前一次性填好所有 GOT 槽，所以 PLT0 永远不会被执行；如果真的走到那里，它会跳到 \code{GOT[2]} 中的 0，现象就是 \code{elr=0} 的异常。
\end{enumerate}

\begin{figure}[htbp]
  \centering
  \includegraphics[width=\linewidth]{echo-dynlink-steps.pdf}
  \caption{运行时动态绑定：\file{/bin/echo} 如何调用到 \file{libc.so} 里的 \code{exit()}}
  \label{fig:echo-dynlink}
\end{figure}

\paragraph{运行时绑定.} 图~\ref{fig:echo-dynlink} 的左列是内核和磁盘上的三个文件，中列是进程的虚拟地址空间（地址向下增长），右列是步骤说明。每条箭头都标有步骤序号：1--3 由内核 \code{exec} 完成，4--6 由 EL0 的 ld-xv6 完成，7--8 是程序本身的运行。各步骤见表~\ref{tab:echo-steps}。

\begin{table}[htbp]
  \centering
  \caption{\code{exit()} 的绑定步骤}
  \label{tab:echo-steps}
  \small
  \begin{tabularx}{\linewidth}{@{}clX@{}}
    \toprule
    步骤 & 执行者 & 内容 \\
    \midrule
    1 & 内核 & 读 \file{/bin/echo} 的程序头，把 LOAD~\#1 装到 \code{0x0}，LOAD~\#2 装到 \code{0x1358} \\
    2 & 内核 & 发现 \code{PT\_INTERP}，把 \file{ld-xv6.so} 装到 \code{PGROUNDUP(0x14b0) = 0x2000}；分配保护页和 4 页栈（\code{0x4000}--\code{0x9000}），压入 argv、envp 和 auxv（\code{AT\_PHDR=0x40}，\code{AT\_BASE=0x2000}，\code{AT\_ENTRY=0x2c8} 等） \\
    3 & 内核 & \code{elr = 0x2000 + 0x750 = 0x2750}（\code{ldso\_start}），\code{eret} 回到 EL0 \\
    4 & ld-xv6 & 用 \code{argv + argc + 2} 找到 auxv，用 \code{AT\_PHDR} 扫描程序头，找到 \code{PT\_DYNAMIC 0x1358}，读出 \code{NEEDED}、\code{SYMTAB}、\code{HASH}、\code{JMPREL} \\
    5 & ld-xv6 & \code{open("/lib/libc.so")}，\code{sbrk} 得到基址 \code{0x9000}，\code{pread} 两个 PT\_LOAD；先处理 libc 自己的 71 个 \code{JUMP\_SLOT} \\
    6 & ld-xv6 & 处理 \code{\_echo} 的 \code{.rela.plt}：\code{sym=2} 是 \code{"exit"}，在主程序中是 UND，于是查 libc 的 \code{.hash} 和 \code{.dynsym}，得到 \code{0x3cc0}，写入 \code{[0x14a0] = 0x9000 + 0x3cc0 = 0xccc0}；write、strlen 同理 \\
    7 & ld-xv6 & 以 \code{(argc, argv)} 调用 \code{AT\_ENTRY = 0x2c8}，即 \code{main} \\
    8 & 程序 & \code{bl exit@plt}（\code{0x2a8}）$\to$ \code{ldr x17, [0x14a0]} 得到 \code{0xccc0} $\to$ \code{br x17} 进入 libc 的 \code{exit} $\to$ \code{svc \#0} \\
    \bottomrule
  \end{tabularx}
\end{table}

关于这张图还有几点说明：

\begin{itemize}
  \item \textbf{libc 基址 \code{0x9000} 的来历。} \file{ld-xv6.so} 的段止于 \code{0x3d20}，向上取页得 \code{0x4000}；保护页和 4 页栈占到 \code{0x9000}，所以 \code{exec} 结束时 \code{sz = 0x9000}，ld-xv6 的第一次 \code{sbrk} 就从这里开始。libc 代码段是 \code{0x9000}--\code{0x122ba}，数据段是 \code{0x132c0}--\code{0x13650}。程序更大时这些值整体上移，例如 \file{/bin/sh} 的 libc 在 \code{0xa000}。
  \item \textbf{内核从不查符号。} 步骤 1--3 只负责按程序头摆放字节、建栈、跳到解释器；符号查找和 GOT 填写全部在 EL0 完成。
  \item \textbf{绑定只发生一次。} 步骤 6 之后，每次库函数调用都只是“\code{bl} 到 PLT，读 GOT，\code{br} 跳转”三步，不再经过 ld-xv6。
  \item \textbf{真机注意 I-cache。} 步骤 5 中 libc 的代码是以数据的形式写入堆页的，在 Cortex-A53 上执行之前需要先做 I-cache 维护，见第~\ref{sec:ldxv6} 节的“问题与建议”。
\end{itemize}

\subsection{进一步：用 \code{mmap} 装载并共享（实验分支）}
\label{sec:ldso}

\file{ld-xv6.so} 让每个进程都有一份私有的 libc。实验分支 \code{wip-cow-mmap-ldso} 在第~\ref{sec:sharedpages} 节的引用计数、写时复制、页缓存和 \code{mmap} 之上，换了一个用 \code{mmap} 装库的解释器 \file{/lib/ld.so}：libc 的代码页全系统只有一份，数据页写时复制。本小节描述这个版本，它尚未合入当前源码树。

\paragraph{链接方式.} 用户程序改为

\begin{makecode}
USER_LD_MODE = -pie -dynamic-linker /lib/ld.so --hash-style=sysv
\end{makecode}

于是每个程序都带上 \code{PT\_INTERP /lib/ld.so} 和 \code{PT\_PHDR}；\code{--hash-style=sysv} 保证有 \code{DT\_HASH} 供符号查找。模式戳改名为 \file{.dynamic-ldso}，旧的二进制会被重新链接。\file{fs.img} 和 ext2 根镜像同时安装 \file{/lib/libc.so} 与 \file{/lib/ld.so}（\code{mkfs} 支持多个 \code{-L}）。

\paragraph{\code{exec()} 的工作.} 主程序照旧装到虚拟地址 0。若有 \code{PT\_INTERP}，读出解释器路径，把它装在主镜像之后（只读页同样取自页缓存，所有进程共用一份 \code{ld.so} 代码），然后\emph{从解释器的入口}开始运行。程序的信息通过辅助向量（auxv）传过去，标签取 Linux 的值：

\begin{txtcode}
x0 = argc   x1 = argv   x2 = auxv（{tag, value} 对，以 AT_NULL 结束）
AT_PHDR    主程序程序头的地址        AT_PHENT/AT_PHNUM  表项大小与个数
AT_PAGESZ  4096                      AT_BASE            ld.so 的装载地址
AT_ENTRY   主程序的入口
\end{txtcode}

Linux 把 argc、argv、envp、auxv 全部放在栈上；这里沿用 xv6 “参数放在寄存器里”的传统，只多传一个 \code{x2}。没有 \code{PT\_INTERP} 的程序仍走上一小节的内核装载器，作为兼容路径。

\paragraph{\code{ld.so} 的工作.} \file{user/ldso.c} 是独立的（不依赖 libc——libc 正是它要装的东西），用原始 \code{svc} 做系统调用，全部符号 hidden，并以 \code{-shared -Bsymbolic -z text} 链接：

\begin{enumerate}
  \item \textbf{自重定位}：根据 \code{AT\_BASE} 处理自己 \code{\_DYNAMIC} 里的 \code{R\_AARCH64\_RELATIVE}。在此之前只能用 PC 相对寻址访问代码和数据。实际构建出的 \code{ld.so} 没有任何重定位项，这一步通常什么也不做；
  \item \textbf{找到程序}：由 \code{AT\_PHDR} 与 \code{PT\_PHDR} 算出主程序基址，读它的 \code{PT\_DYNAMIC}；
  \item \textbf{装载依赖}：对每个 \code{DT\_NEEDED}（不含 \code{/} 的名字到 \file{/lib} 下找），先 \code{mmap} 文件第 0 页读出 ELF 头和程序头，再用一次 \code{mmap(NULL, span, ...)} 把整个镜像范围映射成文件内容（决定了基址），最后对其余段用 \code{MAP\_FIXED} 按各自的地址和权限重新映射，并把 \code{.bss} 清零、不足的部分补匿名页。依赖按广度优先装载，可以有多个库；
  \item \textbf{重定位}：先库后程序，处理 \code{RELATIVE}、\code{ABS64}、\code{GLOB\_DAT}、\code{JUMP\_SLOT}；符号用 SysV 哈希在全局作用域“程序 $\to$ 各库”中查找，所以程序可以覆盖（interpose）libc 的符号；
  \item \textbf{运行}：以 \code{(argc, argv)} 调用 \code{AT\_ENTRY}，用它的返回值调用 \code{exit}。
\end{enumerate}

\paragraph{为什么库不能有文本重定位，程序却可以.} 用户程序是以 \code{-fno-pie} 编译、\code{-pie} 链接的，只读段里的指针表（例如 usertests 的测试表）需要重定位，于是带 \code{DT\_TEXTREL}。程序是 \code{exec()} 装进来的私有可写副本，改写无妨；库的代码页却是页缓存里所有进程共用的只读页，所以 \code{ld.so} 只对库拒绝 \code{DT\_TEXTREL}，libc 本身也以 \code{-z text} 链接来保证这一点。

\paragraph{结果.} 开机后从 \code{/init} 起所有程序都经由 \code{ld.so} 启动。在 QEMU 中（用 gcc 13 构建，libc 代码段为 12 页，可写段从文件第 \code{0xb000} 页开始），\code{vmmap} 显示 \code{init}、\code{sh}、\code{vmmap} 三个进程里：

\begin{itemize}
  \item libc 映射在 \code{0x10\_0000\_1000} 起的 mmap 区；12 个代码页在三个进程里是同一组物理页，每页 “shared by 4 mappings”（页缓存 + 3 个进程）；
  \item 文件第 \code{0xb000} 页既是代码段的尾巴（只读可执行），又是 \code{.data} 的开头（只读 + COW），因此被引用 7 次；
  \item 被重定位写过的 GOT 页已经各自复制成私有的 RW 页——只有真正被写的页才有副本；
  \item \code{ld.so} 的代码页也只有一份物理页。
\end{itemize}

与 Linux 相比还缺：环境变量与 \code{AT\_EXECFN}、\code{dlopen}/\code{dlsym}、延迟绑定、\code{IFUNC}、TLS 与 ASLR；页缓存也只服务 mmap 和库代码，普通 \code{read()} 不经过它。

\subsection{调试动态链接}

出问题时先看是哪一条路径：内核打印 \code{exec: dynamic /bin/x, libc.so base=\ldots} 的是内核内置加载器（没有 \code{PT\_INTERP} 的程序）；不打印却正常运行的走 \code{ld.so}，在 \code{vmmap} 里能看到 libc 位于 \code{0x10\_0000\_xxxx}；\code{DYNAMIC\_USER=0} 的静态程序两者都不经过。主机上常用的检查：

\begin{shcode}
aarch64-elf-readelf -lW user/_sh                # 有没有 INTERP、PHDR；段布局
aarch64-elf-readelf -dW user/libc.so | grep TEXTREL   # 库必须没有
aarch64-elf-readelf -rW user/libc.so            # 重定位类型是否都受支持
aarch64-elf-objdump -s -j .got.plt user/_sh     # GOT 初值
\end{shcode}

表~\ref{tab:dynerr} 列出常见的报错。

\begin{table}[htbp]
  \centering
  \caption{动态链接的常见报错}
  \label{tab:dynerr}
  \small
  \begin{tabularx}{\linewidth}{@{}>{\raggedright\arraybackslash}p{58mm}X@{}}
    \toprule
    信息 & 原因 \\
    \midrule
    \code{exec: \ldots cannot open interpreter /lib/ld.so} & 根文件系统里没装 \code{ld.so}（\code{fs.img} 或 ext2 镜像） \\
    \code{ld.so: cannot open /lib/libc.so} & 缺库，或 \code{DT\_NEEDED} 名字不对 \\
    \code{ld.so: undefined symbol foo} & 程序用了 libc 没有导出的函数 \\
    \code{ld.so: unsupported relocation type} & 出现表~\ref{tab:dynrel} 以外的类型，如 \code{R\_AARCH64\_COPY} \\
    \code{ld.so: text relocations are not supported} & 库里有没用 \code{-fPIC} 编译的对象 \\
    \code{exec: unresolved dynamic symbol} & 同上，发生在内核内置加载器里 \\
    \code{usertrap(): \ldots ec=0x0 \ldots elr=0x0}，随后 \code{init exiting} & 程序是 \code{ET\_DYN}，内核却没有做动态装载：GOT 保持 0，第一次库调用跳到地址 0，执行 ELF 头 \code{\textbackslash x7fELF} 得到未定义指令 \\
    \bottomrule
  \end{tabularx}
\end{table}

最后一条是合并分支时真实遇到的：\code{rpi3-kobject} 的 Makefile 已经把程序链接成动态的，\code{exec.c} 却还是静态装载器，于是 \code{init}（pid 14）在第一次调用 libc 时就跳到了 0。完整的逐步说明见 \file{readme/readme\_dynamic\_libc.md}。

\subsection{POSIX 层}

目标是让只依赖常见 POSIX 接口的 C 源码能直接用 xv6 工具链重新编译。它由三部分组成：

\begin{enumerate}
  \item \file{include/} 提供标准名字的头文件：\code{stdio.h}、\code{stdlib.h}、\code{string.h}、\code{unistd.h}、\code{fcntl.h}、\code{errno.h}、\code{time.h}、\code{pthread.h}、\code{semaphore.h}、\code{sys/types.h}、\code{sys/stat.h}、\code{sys/socket.h}、\code{netinet/in.h}、\code{arpa/inet.h}、\code{sys/epoll.h}；
  \item \file{user/posix.c} 把标准 API 映射到 xv6 已有的能力；
  \item 内核补上无法只靠库函数实现的语义，如 \code{lseek()}，以及 \code{socket()} → \code{bind()} → \code{listen()} → \code{accept()} 的生命周期。
\end{enumerate}

\begin{table}[htbp]
  \centering
  \caption{POSIX 接口到 xv6 的映射}
  \small
  \begin{tabularx}{\linewidth}{@{}lX@{}}
    \toprule
    POSIX 接口 & 实现 \\
    \midrule
    \code{open/read/write/lseek/fstat} & 直接对应系统调用；\code{O\_CREAT} 是 \code{O\_CREATE} 的别名 \\
    \code{getpid/getppid/getcwd/gettid} & 进程与当前工作目录查询 \\
    \code{dup/dup2} & 复制、重定向文件描述符 \\
    \code{wait/waitpid} & 回收僵尸子进程；\code{waitpid} 支持指定直接子进程、任意子进程 \code{-1} 与 \code{WNOHANG} \\
    \code{readlink} & 仅 \code{/sys} 实现（sysfs 符号链接），其它文件系统按 \code{EINVAL} 失败 \\
    \code{pthread\_create/join/exit} & \code{clone}/\code{tjoin}/\code{texit}；\code{pthread\_join} 同时回收用户栈 \\
    \code{pthread\_mutex\_*}、\code{sem\_*} & 内核同步对象；支持 \code{PTHREAD\_MUTEX\_INITIALIZER} 延迟初始化 \\
    \code{clock\_gettime(CLOCK\_MONOTONIC)} & \code{clock\_us()}，读 \reg{CNTVCT\_EL0} \\
    \code{sleep/usleep/nanosleep} & 秒语义；原来按 tick 计的 \code{sleep} 改名为 \code{sleep\_ticks} \\
    \code{socket/bind/listen/accept} & 统一的 socket 文件描述符，可用 \code{read/write/close/epoll} \\
    \bottomrule
  \end{tabularx}
\end{table}

\code{posixtest} 只包含标准头文件，验证文件读写与 \code{lseek}、\code{getcwd}、\code{dup2}、\code{waitpid} 与 \code{getppid}、POSIX 信号量、两个 pthread 并发更新共享变量、静态初始化的互斥锁以及 socket 生命周期。

这仍是第一阶段：系统调用失败只返回 $-1$，没有传回具体错误码，\code{errno} 也还不是线程局部的；\code{waitpid} 尚不支持进程组选择以及 \code{WUNTRACED}/\code{WCONTINUED}，\code{WIFEXITED} 暂时只对应 xv6 的普通退出状态；还缺条件变量、目录流和 TCP \code{connect()}。

\section{在 xv6 里编译 C 程序}

\code{tcc} 是一个面向教学的极小 C 编译器，运行在 xv6 用户态，直接输出 \code{exec()} 能加载的 ELF64 文件，不需要汇编器和链接器：

\begin{shcode}
$ edit hello.c
int main() {
  int answer = 6 * 7;
  printf("hello from xv6 C\n");
  return answer - 42;
}
$ tcc hello.c -o hello
$ hello
hello from xv6 C
\end{shcode}

它支持整数局部变量、四则运算、\code{return}、只含字面量的 \code{printf}/\code{puts} 和注释。生成的代码用 \code{x29} 指向 256 字节的局部变量区，表达式结果放在 \code{x0}，字符串地址用 \code{adr} 指令取得，所以无需重定位；输出和退出直接用 \code{svc \#0} 调用 \code{SYS\_write} 和 \code{SYS\_exit}。文件布局是一个 ELF 头、一个程序头、机器码和字符串池。

一个能在自己身上编译并运行程序的系统，才算是有了一点“自举”的味道。

\section{一个大型用户程序：\code{sshd}}

\code{sshd} 有一千多行，是整个系统里最复杂的用户程序，也是对内核各子系统的一次综合考验。它用到了：

\begin{itemize}
  \item TCP socket 与 \code{accept}——网络协议栈；
  \item 密码学全部在用户态完成：基于 TweetNaCl\cite{tweetnacl} 的 X25519 密钥交换和 Ed25519 主机密钥签名，传输层用 AES-128-CTR 加密、HMAC-SHA256 校验；熵来自内核的硬件随机数发生器（\file{kernel/lib/rng.c}）；
  \item \code{/etc/passwd} 口令认证；
  \item PTY：为每个会话分配一对伪终端，shell 的 fd 0/1/2 指向 PTY 从端，\code{sshd} 读写主端并加密转发；
  \item 会话与控制终端：SSH 会话中的 \code{/dev/tty} 指向该会话的 PTY。
\end{itemize}

调试 \code{sshd} 时发现的问题推动了内核的多处改进：字符回显慢暴露了“网卡每 100\,ms 轮询一次”和“TCP \code{write()} 同步等 ACK”两个瓶颈，促成了 SDIO 中断、workqueue 下半部、NAPI 和 TCP 异步发送队列；\code{ps} 输出跑到串口上，促成了 \code{sys\_ps()} 改为返回缓冲区。

\begin{remark}
  从 \code{\_entry} 的第一条 \code{mrs} 指令，到一个可以通过 Wi-Fi 被 SSH 登录、能在多核上运行实时任务、能用 \code{tcc} 编译程序的系统，中间经过的每一层——串口、页表、进程、文件系统、驱动、协议栈、libc——都写在这本书里了。它们都还很简单，但每一层都在真机上跑过、出过错、被修好过。这也许就是“从裸机到操作系统”最有价值的部分。
\end{remark}
